\documentclass[a4paper,10pt]{article}
\usepackage[utf8]{inputenc}
\usepackage{graphicx}
\usepackage{subfigure}
\usepackage{fancyhdr}
\usepackage{listings}
\usepackage{ gensymb }


%opening

\pagestyle{fancy} %eigener Seitenstil
\fancyhf{} %alle Kopf- und Fußzeilenfelder bereinigen
\fancyhead[L]{SEP Abgabegespräch SS2015/16} %Kopfzeile links
\fancyhead[C]{} %zentrierte Kopfzeile
\fancyhead[R]{\includegraphics[width=1cm]{../../template/tugraz_bw}} %Kopfzeile rechts
\renewcommand{\headrulewidth}{0.4pt} %obere Trennlinie
\fancyfoot[C]{} %Seitennummer
\renewcommand{\footrulewidth}{0.4pt} %untere Trennlinie


\lstset{language=C++, basicstyle=\ttfamily\tiny}  

\begin{document}

\section*{Handymasten}

Schreibe ein Programm, das den nähesten Handymasten des eigenen Providers anzeigt. Dazu soll eine Liste von Handymasten aus einer Datei eingelesen werden. Jede Zeile hat das Format \texttt{<latitude> <longitude> <provider>}. Die aktuelle Position und der eigenen Provider werden als Parameter an das Programm übergeben. 

Zur Berechnung der Entfernung von zwei Koordinaten kann folgende Formel verwendet werden:
\begin{lstlisting}[basicstyle=\footnotesize]
double distance(double lat1, double lon1, double lat2, double lon2) {
  return (lat1 - lat2) * (lat1 - lat2) + (lon1 - lon2) * (lon1 - lon2);
}
\end{lstlisting}

Gibt es keinen Handymasten des eigenen Providers, muss nichts ausgegeben werden. Ist die Anzahl der Parameter falsch, soll ``\texttt{Usage: ./cell-tower <cell-file> <latitude> <longitude> <provider>}'' ausgegeben werden

\vspace{1em}

Beispiel für eine Handymasten-Liste \texttt{cells.txt}:
\begin{lstlisting}[frame=single]
47.065769 15.455070 A1
47.065382 15.445575 Drei
47.059901 15.458664 A1
47.061092 15.444797 A1
47.060555 15.467495 Drei
47.051789 15.466828 Drei
\end{lstlisting}

\subsection*{Beispiel}
\begin{verbatim}
./cell-tower cells.txt 47.056978 15.459423 Drei
Drei (47.0606, 15.4675)

./cell-tower cells.txt 47.056978 15.459423 A1
A1 (47.0599, 15.4587)

./cell-tower cells.txt 47.056978 15.459423 T-Mobile
\end{verbatim}

\subsection*{Aufruf}
\texttt{./cell-tower <cell-tower list> <latitude> <longitude> <provider>}

\subsection*{Bewertung}
\begin{itemize}
 \item Dokumentation und Programmierstil: 1 Punkt
 \item Einhalten der OOP Konzepte: 1 Punkt
 \item Struktur \& Korrektheit: 1 Punkt
\end{itemize}

\newpage


\end{document}
